\section{Reductions}\label{sec:reductions}

\begin{definition}
A bounded operator $T$ on $H$ is \emph{transitive} if its only closed invariant subspaces are
$\{0\}$ and $H$. The invariant subspace problem asks whether a transitive operator exists on a
separable infinite-dimensional $H$. A vector $y$ is \emph{cyclic} if
$\overline{\operatorname{span}}\{T^jy:j\ge0\}=H$.
\end{definition}

To prove there is no transitive operator one may argue by contradiction. This removes every
``\dots or else we already have an invariant subspace'' branch in \paper{}, at the cost that no
concrete operator satisfies the hypotheses. Statements below therefore say explicitly which
consequence of transitivity they use, so that they can be tested on concrete operators.

\begin{lemma}\label{lem:transitive}\stproved{}
Let $T$ be transitive on $H$ with $\dim H=\infty$. Then:
\begin{enumerate}[label=(\roman*)]
\item every non-zero vector is cyclic for $T$;
\item $T^*$ is transitive;
\item if $A$ is bounded and $AT=TA$, then $A=0$ or $A$ is injective with dense range;
\item for every non-zero polynomial $p$, $p(T)$ is injective with dense range;
\item if $g(z)=\sum_jc_jz^j$ has radius of convergence $>\norm{T}$ and $g\not\equiv0$, then
$g(T)$ is injective.
\end{enumerate}
\end{lemma}
\begin{proof}
(i) The closed span of an orbit is a closed invariant subspace containing the vector.
(ii) $M$ is invariant for $T$ iff $M^\perp$ is invariant for $T^*$.
(iii) $\ker A$ and $\overline{AH}$ are closed and $T$-invariant.
(iv) By (iii) each $T-\lambda$ is $0$ or injective with dense range, and $T\ne\lambda$ because
$\dim H>1$ and scalars are not transitive. A product of injective operators with dense range has
both properties.
(v) $g$ is holomorphic on a disc of radius $>\norm T$, so it has finitely many zeros in
$\{|z|\le\norm T\}\supseteq\sigma(T)$. Write $g=pe$ with $p$ a polynomial and $e$ holomorphic and
zero-free on a neighbourhood of $\sigma(T)$. Then $e(T)$ is invertible by the holomorphic
functional calculus and $g(T)=p(T)e(T)$ is injective by (iv).
\end{proof}

\begin{corollary}\label{cor:normalise}\stproved{}
If a transitive operator exists on $H$, then for every $\kappa>0$ there is a transitive $T$ with
$\norm{T}=\kappa$, $T$ injective with dense range, $T$ not surjective and $0\in\sigma(T)$; the same
then holds for $T^*$.
\end{corollary}
\begin{proof}
Pick $\lambda\in\sigma(T_0)$ and put $T=\mu(T_0-\lambda)$ with $\mu>0$ chosen to fix the norm. $T$
has the same invariant subspaces, is injective with dense range by
Lemma~\ref{lem:transitive}(iv), and is not invertible, hence not surjective.
\end{proof}

\noindent\textbf{Standing assumptions.} From now on $T$ is a bounded operator with
$\norm T=\kappa\le\frac12$, injective, with dense range, not surjective. Transitivity is assumed
only where stated.
